% SPDX-License-Identifier: CC-BY-4.0
\documentclass[11pt,letterpaper]{article}

\usepackage[T1]{fontenc}
\usepackage{amsmath,amssymb}
\usepackage{tikz}
\usepackage{pgfplots}
\pgfplotsset{compat=1.18}
\usepackage{algorithm}
\usepackage{algpseudocode}
\usepackage{booktabs}
\usepackage{enumitem}
\usepackage[hidelinks]{hyperref}

\newcommand{\N}{\mathbb{N}}
\newcommand{\Z}{\mathbb{Z}}
\newcommand{\Q}{\mathbb{Q}}
\newcommand{\R}{\mathbb{R}}
\newcommand{\C}{\mathbb{C}}
\newcommand{\Map}[3]{#1 \colon #2 \to #3}
\newcommand{\Mod}[1]{#1\text{-}\mathrm{Mod}}
\newcommand{\GroupHom}{\mathrm{Hom}_{\mathbf{Grp}}}
\newcommand{\RingHom}{\mathrm{Hom}_{\mathbf{Ring}}}

\newtheorem{theorem}{Theorem}
\newtheorem{lemma}[theorem]{Lemma}
\newtheorem{corollary}[theorem]{Corollary}
\newtheorem{definition}[theorem]{Definition}
\newtheorem{remark}[theorem]{Remark}
\newenvironment{proof}{\noindent\textit{Proof.}}{\hfill\(\square\)}

\title{Generalized Collatz Library}
\author{}
\date{\today}

\begin{document}
\maketitle

\begin{abstract}
We distinguish exact Gaussian and Eisenstein lattice maps, residue-wise
quotient variants, affine permutations on finite quadratic rings, and
selected holomorphic polynomial models. The raw \((3,c,2)\) nonreal escape theorem has a finite dyadic descent proof. A broader scalar theorem covers arbitrary real integral offsets and signed multiplicatively independent magnitudes, including an escaping orbit that divides infinitely often. The quotient variants have nonreal cycles. Four polynomial models preserve the same selected real orbit values but have different exact cycle multipliers. Their differing multipliers show that finite arithmetic orbit data do
not determine complex stability. Their exact escape radius is proved, while the plotted pixels remain finite-time floating-point experiments. The numerical census is bounded and all conventions, clocks, fences, and unresolved classifications are retained.
\end{abstract}

\section{Four different dynamical objects}

For \(R=\Z[\theta]\), with \(\theta=i\) or \(\omega=e^{2\pi i/3}\), the raw map is
\[
T(z)=\begin{cases}z/d,&d\mid z\text{ in }R,\\qz+c,&d\nmid z\text{ in }R.\end{cases}
\]
It is not a map using literal field divisibility: over a field, every nonzero divisor divides every point and only the quotient branch survives. The state ring must therefore be retained.

For scalar divisor \(2\), the separately defined residue quotient maps are
\[
J_\pm(z)=\begin{cases}z/2,&r(z)=0,\\(3z\pm r(z))/2,&r(z)\ne0,\end{cases}
\qquad r(z)=(x\bmod2)+(y\bmod2)\theta.
\]
Their numerators are integral. The Gaussian plus rule is the Joseph convention printed in Kucinski~\cite{Kucinski00}; the Eisenstein and minus variants are explicitly specified here, not attributed to that paper. These maps have their own ticks, not raw fixed-offset clocks. Finally, \(z\mapsto\theta z+1\) on a finite quadratic quotient is a genuine affine permutation. Raw halving does not automatically descend to that quotient. Polynomial iteration on \(\C\) is a fourth object, with a separately specified holomorphic function.

Geometric study of higher-rank residue-wise Collatz mappings has prior literature~\cite{Kionke17}. Its conventions and cone hypotheses must be checked before importing a divergence statement; no such theorem is used in the elementary raw descent below. Real analytic Collatz extensions also have a separate literature~\cite{Chamberland96}; neither real interpolation nor agreement at integer points selects a unique complex polynomial.

\section{Exact lattice arithmetic and finite survey}

Write \(\theta^2=\tau\theta-\nu\). For \(d=s+t\theta\), multiplication by \(d\) has coefficient matrix
\[
\begin{pmatrix}s&-\nu t\\t&s+\tau t\end{pmatrix},
\qquad N(d)=s^2+\tau st+\nu t^2.
\]
For nonzero \(d\) in the chosen imaginary quadratic order, the determinant is positive. The adjugate gives
\[
d\mid x+y\theta
\quad\Longleftrightarrow\quad
N(d)\mid(s+\tau t)x+\nu ty
\quad\text{and}\quad N(d)\mid-tx+sy.
\]
Both divisibilities are necessary and sufficient because the two quotient coefficients are those numerators divided by \(N(d)\). This requires no element unique-factorization assumption. The companion quadratic arithmetic development supplies the coordinate formalization; Go tests compare against independent rational matrix elimination.

For the scalar raw \((3,c,2)\) map with odd real \(c\), write any nonzero second coefficient as \(2^r u\), \(u\) odd. An even first coefficient permits immediate division. An odd first coefficient takes one affine tick to become even and then divides. Each one- or two-tick block removes one factor of two from the second coefficient. After at most \(2r\) ticks it is odd and then equals \(3^j y_t\). Consequently every nonreal seed escapes along its full sequence. The companion \texttt{ComplexLattice.lean} checks the rank descent and the quantified transverse-height conclusion; the written lattice proof supplies the complex embedding bound.

In contrast, \(J_+\) has \(1+\theta\to2+2\theta\to1+\theta\), while \(J_-\) fixes \(1+\theta\). Those nonreal quotient cycles are explicit negative controls for conflating the conventions. Figure~\ref{fig:lattices} juxtaposes the two rules.

\begin{figure}[htbp]
\centering
\includegraphics[width=.47\textwidth]{../other/GaussianMinus.pdf}
\hfill
\includegraphics[width=.47\textwidth]{../other/GaussianJosephPlus.pdf}
\caption{Left: raw Gaussian \((3,-1,2)\), with proved off-axis escape. Right: the Gaussian plus residue-quotient convention, with nonreal reached cycles. Each cell represents one integral seed, not a continuous orbit. Navy denotes a theorem-backed escaping tail, orange only a coordinate-fence exit, gray an unresolved prefix, and other colors reached exact cycles. The six-color cycle hash is not injective; the marked census identifies each cycle.}
\label{fig:lattices}
\end{figure}

The exact workload is thirteen families, each with \(97^2=9,409\) seeds in the coefficient box \([-48,48]^2\), at most \(256\) ticks, and coordinate fence \(2^{26}\). The total is \(122,317\) seeds. For Gaussian coordinates plot \((x,y)\); for Eisenstein coordinates plot \((x-y/2,\sqrt3y/2)\). Repeated scalar-law counts in these embeddings are not independent convergence evidence.

The four raw \((3,\pm1,2)\) rows each have \(97\) exact-cycle seeds and \(9,312\) proved off-axis escapes. Each of the four quotient rows has \(745\) exact-cycle seeds, \(8,510\) fence exits, and \(154\) unresolved prefixes, reaching \(45\) distinct cycles from the window. For the genuinely complex raw parameters, cycle seeds are \(13\) for Gaussian \((2,1,1+i)\), \(1\) for Gaussian \((2i,1+i,2)\), \(1\) for Eisenstein \((2\omega,1,2)\), \(8\) for Eisenstein \((1-\omega,1,1-\omega)\), and \(1\) for Gaussian \((2+i,0,2-i)\). The remaining seeds in those five rows cross the fence. They are not called universally escaping from this finite test. Complete per-family counts, cycle representatives, and primitive periods are retained in \texttt{DynamicsData.json}.

\section{Independent scalar families and a dependent control}

The general scalar result and actual ring-to-complex recovery are proved in
\texttt{QCD.\allowbreak Complex\allowbreak Lattice.\allowbreak 001/\allowbreak tex/\allowbreak Scalar\allowbreak Escape.tex} and
\texttt{QCD.\allowbreak Complex\allowbreak Lattice.\allowbreak 001/\allowbreak lean/\allowbreak Scalar\allowbreak Escape.lean}
\cite{LibraryScalar26}. For
\(q,c,d\in\Z\), \(|q|,|d|\ge2\), it assumes
\[
|q|^a\ne|d|^b\qquad(a,b\in\N_{>0})
\]
and proves full-sequence complex-norm escape for every nonreal seed on any
of the specified imaginary quadratic coordinate orders. It does not
classify the real integer axis.

Here is the arithmetic hinge. If \(a_n,b_n\) count affine and division
ticks, legal exact division and scalar multiplication give
\[
a_n+b_n=n,\qquad d^{b_n}y_n=q^{a_n}y_0.
\]
A repeated nonzero transverse magnitude at two times would, on that
actual intervening segment, imply \(|q|^a=|d|^b\) with \(a,b>0\), a
contradiction. The resulting injective sequence of nonnegative integer
magnitudes visits every bounded set only finitely often. Its divergence
recovers complex-norm escape through the positive imaginary basis scale.
This argument retains the two-coordinate guard; it is not a scalar
one-coordinate raw quotient.

For \(12=2^2\cdot3\) and \(18=2\cdot3^2\), the prime-exponent minor
is \(2\cdot2-1\cdot1=3\ne0\). The raw \((12,0,18)\) seed \(\theta\)
has the exact three-tick blocks
\[
8^j\theta\to12\cdot8^j\theta\to144\cdot8^j\theta\to8^{j+1}\theta.
\]
It divides forever. An eventual division-free tail is therefore not a
valid prerequisite for the general result. In contrast, the dependent
pair \((4,8)\) has the primitive five-cycle
\[
\theta\to4\theta\to16\theta\to2\theta\to8\theta\to\theta.
\]
Both examples and the independence minor are checked in Lean.

Six additional scalar families each use the same \(9,409\) seeds,
\(256\)-tick horizon and \(2^{26}\) census fence: \(56,454\) added seeds,
\(178,771\) combined lattice seeds. The JSON stores these separately as
\texttt{ScalarLattices}. Gaussian \((12,0,18)\) has \(9,408\)
theorem-backed escapes and one exact-cycle seed, including \(96\)
nonzero real-axis seeds certified by the homogeneous theorem.
The Gaussian \((3,2,2)\), \((2,7,3)\), signed \((-2,-3,-3)\),
and Eisenstein \((12,5,18)\) rows each have \(9,312\) theorem-backed
off-axis escapes. Their respective real-axis exact-cycle/fence-exit
counts are \((7,90)\), \((19,78)\), \((5,92)\), and \((1,96)\). The Gaussian \((4,0,8)\)
row has \(9,409\) cycle seeds reaching \(7,009\) marked cycles in this
window. None of these rows has an unresolved prefix. A real-axis fence
exit is not certified escape, and the dependent row is not an infinite
all-seed cycle enumeration.

\begin{figure}[htbp]
\centering
\includegraphics[width=.47\textwidth]{../other/GaussianCommonSupport.pdf}
\hfill
\includegraphics[width=.47\textwidth]{../other/GaussianDependentScalar.pdf}
\caption{Independent shared-support scalar parameters \((12,0,18)\) on the
left, and dependent \((4,0,8)\) on the right. Navy denotes theorem-backed
escape; for \((12,0,18)\), this includes every nonzero real-axis seed.
Other colors identify exact reached cycles by the same noninjective
six-color hash.}
\end{figure}

Coefficient matching between two quadratic orders is a bijective
same-clock conjugacy for fixed scalar parameters, but need not preserve
multiplication: matching \(i\) to \(\omega\) sends \(i^2\) to \(-1\),
whereas \(\omega^2=-1-\omega\). The companion formal proof separates
the dynamical isomorphism from ring and ideal arithmetic.

\section{Concrete cohomology, fields, and ramification}

On an oriented nonempty cycle, the ordinary cochain differential is \((\delta f)_j=f_{j+1}-f_j\). Prefix sums prove that a cochain is exact exactly when its total period is zero. Thus \(H^1\cong K\) over any field, regardless of cycle length. This is checked as an actual finite-cochain quotient \texttt{LinearEquiv} in \texttt{ComplexCohomology.lean}; unmarked cohomology does not recover the length or guards.

For an affine local system \(u_{j+1}=\lambda_j u_j+a_j\), with each \(\lambda_j\ne0\), ordered composition gives
\[
P=\prod_j\lambda_j,
\qquad B=\sum_j\left(\prod_{k>j}\lambda_k\right)a_j,
\qquad (1-P)u_0=B.
\]
For \(P\ne1\), every twisted cochain has a unique periodic potential and \(H^0_\lambda=H^1_\lambda=0\). For \(P=1\), exactness is \(B=0\), and the kernel and cokernel are each isomorphic to \(K\). The companion cohomology proof reconstructs one common potential and states the remaining lattice integrality and branch-guard obligations explicitly.

The affine permutations on \(\Z[i]/(2)\), \(\Z[i]/(3)\), \(\Z[\omega]/(2)\), and \(\Z[\omega]/(3)\) have cycle-length lists \((4)\), \((1,4,4)\), \((1,3)\), and \((3,3,3)\). Their incidence ranks are \(3,6,2,6\) over each independent coefficient field \(\mathbb F_2,\mathbb F_3,\mathbb F_5\). The middle two state rings are \(\mathbb F_9\) and \(\mathbb F_4\); the first and last have nonzero nilpotents and are not fields. Thus finite field/Galois data, state rings, and cochain fields are not interchangeable. Tests exhaust every nonzero-multiplier word of lengths one through three over the three cochain fields against direct enumeration of potentials.

\subsection{Explicit Galois fields and their useful symmetry}

% Root-free quadratic defining polynomials justify genuine field structures.
\begin{theorem}[Quadratic Frobenius and complete automorphism groups]
Let \(K_9=\mathbb F_3[\bar i]\), with \(\bar i^2=-1\), and \(K_4=\mathbb F_2[\bar\omega]\), with \(\bar\omega^2+\bar\omega+1=0\). These have \(9\) and \(4\) elements. Their nontrivial coefficient-field automorphisms are
\begin{align*}
\sigma_9(x+y\bar i)&=x-y\bar i=(x+y\bar i)^3,\\
\sigma_4(x+y\bar\omega)&=(x+y)+y\bar\omega=(x+y\bar\omega)^2.
\end{align*}
In each case every coefficient-field algebra automorphism is identity or \(\sigma\); these are distinct, and \(\sigma^2=1\). The fixed field consists exactly of the elements with \(y=0\).
\end{theorem}
\begin{proof}
The polynomial \(X^2+1\) has no root in \(\mathbb F_3\), whose squares are \(0,1\). The polynomial \(X^2+X+1\) has no root in \(\mathbb F_2\), by evaluating at \(0,1\). Thus the corresponding quadratic algebras are fields. Every element has a unique two-coefficient representation, giving the cardinalities.

In characteristic three, the binomial cross terms vanish, \(x^3=x,y^3=y\), and \(\bar i^3=-\bar i\). In characteristic two, squaring is additive and \(\bar\omega^2=1+\bar\omega\). These give the displayed coordinate formulas, which square to identity and fix exactly the points with \(y=0\). They are not identity on the generator.

The conjugate generators are distinct roots of the defining quadratic, so each polynomial splits with two distinct roots in the generated field. A coefficient-field algebra homomorphism is determined by its generator image: it sends \(x+y\alpha\) to \(x+yf(\alpha)\). The defining polynomial requires this image to be one of its two roots. The two exhibited automorphisms therefore exhaust the group, which is the cyclic group of order two. The Lean file \texttt{Complex\allowbreak Fields.lean} checks the root choices, cardinalities, Frobenius formulas and automorphism exhaustion.
\end{proof}

% Both parameters and every tick are retained by the actual algebra automorphism.
\begin{corollary}[Affine orbit transport]
For \(T_{q,c}(z)=qz+c\), quadratic conjugation satisfies
\[
\sigma(T_{q,c}^{\circ j}(z))
=T_{\sigma(q),\sigma(c)}^{\circ j}(\sigma(z))
\]
at every nonnegative clock \(j\). The inverse reconstruction is the same formula because \(\sigma^2=1\).
\end{corollary}
\begin{proof}
The case \(j=0\) is immediate. For the next tick, preservation of multiplication and addition gives
\[
\sigma(qT_{q,c}^{\circ j}(z)+c)
=\sigma(q)\sigma(T_{q,c}^{\circ j}(z))+\sigma(c).
\]
Insert the induction hypothesis. The original state and both parameters recover by applying the involution again. Thus conjugate multipliers with \(c=1\) have isomorphic marked cycle graphs and identical incidence ranks. Lean checks the full orbit identity as \texttt{ConjugationAffineOrbit}. This finite-model symmetry does not remove lattice branch guards.
\end{proof}

The opposite-characteristic controls are not fields: \(1+\bar i\) modulo \(2\) and \(1-\bar\omega\) modulo \(3\) are nonzero with square zero. The Go census exhausts all elements and finds Frobenius image sizes \(2,9,4,3\) and counts of invertible nonzero elements \(2,8,3,6\), in the ring order used above. Frobenius is not injective on the two ramified rings.

\section{Polynomial Julia countermodels}

For a polynomial \(p\) of degree at least two, define its filled Julia set and Julia set by
\[
K_p=\{z\in\C:\sup_j|p^{\circ j}(z)|<\infty\},
\qquad J_p=\partial K_p.
\]
The four models below have these genuine holomorphic dynamical objects. They preserve only specified finite arithmetic nodes, not the entire integer or complex lattice raw rule.

Choose
\begin{align*}
p_-(z)&=-\frac32z^2+\frac72z,&
\widetilde p_-(z)&=\frac12z^3-3z^2+\frac92z,\\
p_+(z)&=\frac76z^3-7z^2+\frac{59}6z,&
\widetilde p_+(z)&=\frac29z^4-\frac7{18}z^3-\frac{35}9z^2+\frac{145}{18}z.
\end{align*}

% Direct node checks establish agreement only on the stated finite cycles.
\begin{lemma}[Finite node agreement only]
Both minus models fix \(0\) and preserve \(1\to2\to1\). Both plus models fix \(0\) and preserve \(1\to4\to2\to1\). They are not raw integer extensions: already at \(3\), the first three models above take values \(-3,0,-2\), whereas the raw minus or plus value is \(8\) or \(10\), respectively. The fourth gives \(-10/3\), also not \(10\).
\end{lemma}
\begin{proof}
Substitution gives the stated finite values. The perturbations are
\[
\widetilde p_--p_- =\tfrac12 z(z-1)(z-2),
\qquad \widetilde p_+-p_+ =\tfrac29 z(z-1)(z-2)(z-4).
\]
These vanish at the corresponding finite nodes. Substituting \(z=3\) gives the failure values, including \(p_+(3)=-2\) and the fourth value \(-2+(2/9)(3\cdot2\cdot1\cdot(-1))=-10/3\). No all-integer agreement follows from the finite-node identities. The companion Lean theorem checks all four failure values and the finite node identities.
\end{proof}

% Cycle multipliers demonstrate different complex stability despite identical node values.
\begin{theorem}[MinusMultipliers and PlusMultipliers]
The minus two-cycle has derivative-product multipliers \(-5/4\) for \(p_-\) and \(0\) for \(\widetilde p_-\). The plus three-cycle has multipliers \(1475/54\) for \(p_+\) and \(0\) for \(\widetilde p_+\).
\end{theorem}
\begin{proof}
For \(p_-\), derivatives at \(1,2\) are \(1/2,-5/2\), whose product is \(-5/4\). For \(\widetilde p_-\), the derivative at \(1\) is \(3/2-6+9/2=0\), so the product is zero. For \(p_+\), the derivative \((7/2)z^2-14z+59/6\) has values \(-2/3,59/6,-25/6\) at \(1,4,2\); their product is \(1475/54\). The derivative of the plus perturbation at \(1\) is \((2/9)(1\cdot(-1)\cdot(-3))=2/3\), cancelling \(p_+'(1)=-2/3\). Hence the perturbed product is zero. All derivative-product values and their identification with actual complex derivatives are checked in Lean.
\end{proof}

If \(g=p^{\circ n}\) is the return map at a cycle point \(z_0\), the complex chain rule makes \(g'(z_0)\) this product. Differentiability gives \(g(z_0+h)-z_0=\mu h+o(|h|)\). If \(|\mu|>1\), choosing the remainder smaller than \((|\mu|-1)|h|/2\) gives local expansion. If \(\mu=0\), choosing it smaller than \(|h|/2\) gives contraction for sufficiently small \(h\). Thus the first models have repelling preserved cycles and the perturbed models have superattracting preserved cycles. This local stability deduction is written here; the compiled polynomial file checks derivatives and exact multipliers, not a general Julia classification theorem.

Figures~\ref{fig:minus} and \ref{fig:plus} display the resulting contrasting escape portraits. Agreement on the arithmetic cycle values cannot determine the complex stability or basin geometry of a chosen interpolation. This is a refuted transfer route, not a proof that lattice behavior is read off a Julia image.

\begin{figure}[htbp]
\centering
\includegraphics[width=.47\textwidth]{../other/MinusQuadratic.pdf}
\hfill
\includegraphics[width=.47\textwidth]{../other/MinusCritical.pdf}
\caption{Finite-time portraits of the two minus models. Both preserve the same two-cycle; its multiplier changes from \(-5/4\) to \(0\). Colors record the first numerical tick exceeding radius \(8\); dark pixels only survived the declared horizon.}
\label{fig:minus}
\end{figure}

\begin{figure}[htbp]
\centering
\includegraphics[width=.47\textwidth]{../other/PlusCubic.pdf}
\hfill
\includegraphics[width=.47\textwidth]{../other/PlusCritical.pdf}
\caption{Finite-time portraits of the two plus models. Both preserve \(1\to4\to2\to1\); the exact multiplier changes from \(1475/54\) to \(0\). These are Julia countermodels for finite-node inference, not all-lattice Collatz extensions.}
\label{fig:plus}
\end{figure}

\section{Exact escape radius and numerical boundary}

% The radius estimate uses the polynomial coefficients rather than sampled angles.
\begin{theorem}[Growth certificates and OutsideEightPowers]
Each of the four models satisfies
\[
|z|\ge8\ \Longrightarrow\ |p(z)|\ge2|z|.
\]
Consequently \(|p^{\circ j}(z)|\ge2^j|z|\) for every \(j\ge0\) if \(|z|\ge8\).
\end{theorem}
\begin{proof}
Write \(R=|z|\ge8\). The triangle inequality applied against the leading term gives the following lower bounds divided by \(R>0\):
\[
\tfrac32R-\tfrac72,\quad
\tfrac12R^2-3R-\tfrac92,\quad
\tfrac76R^2-7R-\tfrac{59}6,\quad
\tfrac1{18}(4R^3-7R^2-70R-145).
\]
The first is at least \(17/2\). The second equals \(\tfrac12R(R-6)-9/2\), at least \(R-9/2\ge7/2\). The third equals \(\tfrac76R(R-6)-59/6\), at least \((7/3)R-59/6\ge53/6\). For the fourth, \(4R-7\ge25\), so its numerator is at least \(25R^2-70R-145\ge130R-145\ge895\). Every displayed lower bound is greater than \(2\). This proves the radius estimate for all angles, not merely sampled ones. The successor stays outside radius \(8\). Induction applies the same estimate at every iterate and gives the power bound. The estimates and iterative recovery are checked in Lean.
\end{proof}

The Go portraits use a \(768\times768\) mesh of sample centers for each view, \(256\) iterations, and threshold \(|z|>8\). Four global views and two local zooms of the same four polynomials give \(3,538,944\) samples total. A separate comparison of \(128\) and \(256\) ticks at \(65,536\) shared centers in each view finds no additional escapes in these finite comparisons. The global plus panels do not resolve the small basin: the two added views use \([0.994,1.006]+i[-0.006,0.006]\) around the preserved point \(1\), without changing coefficients. A fixed-mesh numerical control observes zero finite-horizon survivors for the repelling zoom and a nonempty surviving region for the critical zoom; this does not certify basin membership. Horner iteration uses \texttt{complex128}. Its roundoff is not bounded by interval arithmetic; the exact escape theorem therefore does not certify numerical membership near a Julia boundary. A dark pixel is an unresolved finite-time observation, not a bounded-orbit theorem or a dimension calculation.

\section{Reproduction, visual quality, and remaining obligations}

The generator writes twenty-five vector SVGs, twenty-five vector PDFs, twenty-five PNG previews, and finite data. Adjacent equal-valued samples are coalesced without changing their colors. SVGs contain outlined glyphs and vector cells, not embedded raster images. Each preview is \(1500\times1500\), at 300 dpi for a \(5\) inch square. Latin Modern outlines use identical metrics for measuring and drawing; title, axis, and tick sizes are explicitly 16, 14, and 12 points. The label renderer recognizes only its ten declared expressions and rejects unsupported math. The PDFs are ignored by Git; Markdown displays the vector SVGs. Polynomial PNG previews use a single raster tile to avoid seams between adjacent cells. Contrast-backed legends identify the displayed categories. Optional annotations mark an exact lattice cycle or the finite-horizon survivor-mask edge, explicitly not an exact Julia boundary.

Go tests verify exact signed division independently, finite dyadic entry budgets, raw-versus-quotient and field-versus-ring controls, finite incidence ranks, exhaustive short local-system words, selected polynomial values, and warmed zero-allocation orbit diagnosis. Polynomial sampled-angle tests are implementation controls, not the proof of the radius. The ordinary cochain quotient, affine reconstruction, raw off-axis height theorem, and polynomial identities have matching Lean modules.

Established are the stated restricted theorem, explicit quotient cycles, field-guard obstruction, finite graph cohomology calculations, and noncanonical polynomial countermodels. Open are real-axis universal targets, other complex raw families, exhaustive Joseph/Eisenstein infinite-cycle classification, and a justified Julia-to-arithmetic transfer. An informative next step is admissible branch-word enumeration with integral recovery and retained cycle marks. Unmarked cohomology or visual resemblance alone cannot replace those obligations.

\bibliographystyle{amsalpha}
\bibliography{ComplexDynamics}

\end{document}
